% NeurIPS submission using official public style package.
% CHANGELOG (v2):
% - Added Appendix F (Propositions F.1--F.5): broadband decomposition, finite-(b,L) residual bounds, resonance regularization, off-diagonal decay, and quantization bounds.
% - Strengthened Eq. (10) motivation: explicit diagonal ridge + off-diagonal Brownian decomposition and regularization rationale.
% - Added a compact continuous formulation subsection for anchored-sigmoid with parameter interpretation.
% - Removed internal artifact identifiers and local file/path leakage from the manuscript body and appendices.
% - Rewrote Appendix H into formal hyperparameter-configuration tables while preserving anonymous controlled-protocol reproducibility.
% CHANGELOG (v4):
% - Reframed narrative: principled tradeoff analysis rather than leaderboard claims.
% - Added cross-model Qwen-2.5-7B dual-seed (42/1337) 400-step LongBench-21 validation.
% - Added Section 6 "Cross-Model Tradeoff Analysis" with task-family decomposition.
% - Strengthened Limitations with honest tradeoff characterization.
% - Updated abstract and conclusion to reflect theory-first tradeoff framing.
\documentclass{article}

\begin{filecontents*}{neurips_2025.sty}
% partial rewrite of the LaTeX2e package for submissions to the
% Conference on Neural Information Processing Systems (NeurIPS):
%
% - uses more LaTeX conventions
% - line numbers at submission time replaced with aligned numbers from
%   lineno package
% - \nipsfinalcopy replaced with [final] package option
% - automatically loads times package for authors
% - loads natbib automatically; this can be suppressed with the
%   [nonatbib] package option
% - adds foot line to first page identifying the conference
% - adds preprint option for submission to e.g. arXiv
% - conference acronym modified
% - update foot line to display the track name
%
% Roman Garnett (garnett@wustl.edu) and the many authors of
% nips15submit_e.sty, including MK and drstrip@sandia
%
% last revision: April 2025

\NeedsTeXFormat{LaTeX2e}
\ProvidesPackage{neurips_2025}[2025/05/01 NeurIPS 2025 submission/camera-ready style file]

% declare final option, which creates camera-ready copy
\newif\if@neuripsfinal\@neuripsfinalfalse
\DeclareOption{final}{
  \@neuripsfinaltrue
  \@anonymousfalse
}

% declare nonatbib option, which does not load natbib in case of
% package clash (users can pass options to natbib via
% \PassOptionsToPackage)
\newif\if@natbib\@natbibtrue
\DeclareOption{nonatbib}{
  \@natbibfalse
}

% declare preprint option, which creates a preprint version ready for
% upload to, e.g., arXiv
\newif\if@preprint\@preprintfalse
\DeclareOption{preprint}{
  \@preprinttrue
  \@anonymousfalse
}

% determine the track of the paper in camera-ready mode
\newif\if@main\@maintrue
\DeclareOption{main}{
  \@maintrue
  \newcommand{\@trackname}{\@neuripsordinal\ Conference on Neural Information Processing Systems (NeurIPS \@neuripsyear).}
}
\newif\if@position\@positionfalse
\DeclareOption{position}{
  \@positiontrue
  \newcommand{\@trackname}{\@neuripsordinal\ Conference on Neural Information Processing Systems (NeurIPS \@neuripsyear) Position Paper Track.}
}
\newif\if@dandb\@dandbfalse
\DeclareOption{dandb}{
  \@dandbtrue
  \@anonymousfalse
  \newcommand{\@trackname}{\@neuripsordinal\ Conference on Neural Information Processing Systems (NeurIPS \@neuripsyear) Track on Datasets and Benchmarks.}
}
\newif\if@creativeai\@creativeaifalse
\DeclareOption{creativeai}{
  \@creativeaitrue
  \@anonymousfalse
  \newcommand{\@trackname}{\@neuripsordinal\ Conference on Neural Information Processing Systems (NeurIPS \@neuripsyear) Creative AI Track.}
}

% For anonymous or non-anonymous
\newif\if@anonymous\@anonymoustrue

% For workshop papers
\newcommand{\@workshoptitle}{}
\newcommand{\workshoptitle}[1]{\renewcommand{\@workshoptitle}{#1}}

\newif\if@workshop\@workshopfalse
\DeclareOption{sglblindworkshop}{
  \@workshoptrue
  \@anonymousfalse
  \newcommand{\@trackname}{\@neuripsordinal\ Conference on Neural Information Processing Systems (NeurIPS \@neuripsyear) Workshop: \@workshoptitle.}
}
\DeclareOption{dblblindworkshop}{
  \@workshoptrue
  \newcommand{\@trackname}{\@neuripsordinal\ Conference on Neural Information Processing Systems (NeurIPS \@neuripsyear) Workshop: \@workshoptitle.}
}

\ProcessOptions\relax

% fonts
\renewcommand{\rmdefault}{ptm}
\renewcommand{\sfdefault}{phv}

% change this every year for notice string at bottom
\newcommand{\@neuripsordinal}{40th}
\newcommand{\@neuripsyear}{2026}
\newcommand{\@neuripslocation}{Sydney, Australia}

% acknowledgments
\usepackage{environ}
\newcommand{\acksection}{\section*{Acknowledgments and Disclosure of Funding}}
\NewEnviron{ack}{%
  \acksection
  \BODY
}


% load natbib unless told otherwise
\if@natbib
  \RequirePackage{natbib}
\fi

% set page geometry
\usepackage[verbose=true,letterpaper]{geometry}
\AtBeginDocument{
  \newgeometry{
    textheight=9in,
    textwidth=5.5in,
    top=1in,
    headheight=12pt,
    headsep=25pt,
    footskip=30pt
  }
  \@ifpackageloaded{fullpage}
    {\PackageWarning{neurips_2025}{fullpage package not allowed! Overwriting formatting.}}
    {}
}

\widowpenalty=10000
\clubpenalty=10000
\flushbottom
\sloppy


% font sizes with reduced leading
\renewcommand{\normalsize}{%
  \@setfontsize\normalsize\@xpt\@xipt
  \abovedisplayskip      7\p@ \@plus 2\p@ \@minus 5\p@
  \abovedisplayshortskip \z@ \@plus 3\p@
  \belowdisplayskip      \abovedisplayskip
  \belowdisplayshortskip 4\p@ \@plus 3\p@ \@minus 3\p@
}
\normalsize
\renewcommand{\small}{%
  \@setfontsize\small\@ixpt\@xpt
  \abovedisplayskip      6\p@ \@plus 1.5\p@ \@minus 4\p@
  \abovedisplayshortskip \z@  \@plus 2\p@
  \belowdisplayskip      \abovedisplayskip
  \belowdisplayshortskip 3\p@ \@plus 2\p@   \@minus 2\p@
}
\renewcommand{\footnotesize}{\@setfontsize\footnotesize\@ixpt\@xpt}
\renewcommand{\scriptsize}{\@setfontsize\scriptsize\@viipt\@viiipt}
\renewcommand{\tiny}{\@setfontsize\tiny\@vipt\@viipt}
\renewcommand{\large}{\@setfontsize\large\@xiipt{14}}
\renewcommand{\Large}{\@setfontsize\Large\@xivpt{16}}
\renewcommand{\LARGE}{\@setfontsize\LARGE\@xviipt{20}}
\renewcommand{\huge}{\@setfontsize\huge\@xxpt{23}}
\renewcommand{\Huge}{\@setfontsize\Huge\@xxvpt{28}}

% sections with less space
\providecommand{\section}{}
\renewcommand{\section}{%
  \@startsection{section}{1}{\z@}%
                {-2.0ex \@plus -0.5ex \@minus -0.2ex}%
                { 1.5ex \@plus  0.3ex \@minus  0.2ex}%
                {\large\bf\raggedright}%
}
\providecommand{\subsection}{}
\renewcommand{\subsection}{%
  \@startsection{subsection}{2}{\z@}%
                {-1.8ex \@plus -0.5ex \@minus -0.2ex}%
                { 0.8ex \@plus  0.2ex}%
                {\normalsize\bf\raggedright}%
}
\providecommand{\subsubsection}{}
\renewcommand{\subsubsection}{%
  \@startsection{subsubsection}{3}{\z@}%
                {-1.5ex \@plus -0.5ex \@minus -0.2ex}%
                { 0.5ex \@plus  0.2ex}%
                {\normalsize\bf\raggedright}%
}
\providecommand{\paragraph}{}
\renewcommand{\paragraph}{%
  \@startsection{paragraph}{4}{\z@}%
                {1.5ex \@plus 0.5ex \@minus 0.2ex}%
                {-1em}%
                {\normalsize\bf}%
}
\providecommand{\subparagraph}{}
\renewcommand{\subparagraph}{%
  \@startsection{subparagraph}{5}{\z@}%
                {1.5ex \@plus 0.5ex \@minus 0.2ex}%
                {-1em}%
                {\normalsize\bf}%
}
\providecommand{\subsubsubsection}{}
\renewcommand{\subsubsubsection}{%
  \vskip5pt{\noindent\normalsize\rm\raggedright}%
}

% float placement
\renewcommand{\topfraction      }{0.85}
\renewcommand{\bottomfraction   }{0.4}
\renewcommand{\textfraction     }{0.1}
\renewcommand{\floatpagefraction}{0.7}

\newlength{\@neuripsabovecaptionskip}\setlength{\@neuripsabovecaptionskip}{7\p@}
\newlength{\@neuripsbelowcaptionskip}\setlength{\@neuripsbelowcaptionskip}{\z@}

\setlength{\abovecaptionskip}{\@neuripsabovecaptionskip}
\setlength{\belowcaptionskip}{\@neuripsbelowcaptionskip}

% swap above/belowcaptionskip lengths for tables
\renewenvironment{table}
  {\setlength{\abovecaptionskip}{\@neuripsbelowcaptionskip}%
   \setlength{\belowcaptionskip}{\@neuripsabovecaptionskip}%
   \@float{table}}
  {\end@float}

% footnote formatting
\setlength{\footnotesep }{6.65\p@}
\setlength{\skip\footins}{9\p@ \@plus 4\p@ \@minus 2\p@}
\renewcommand{\footnoterule}{\kern-3\p@ \hrule width 12pc \kern 2.6\p@}
\setcounter{footnote}{0}

% paragraph formatting
\setlength{\parindent}{\z@}
\setlength{\parskip  }{5.5\p@}

% list formatting
\setlength{\topsep       }{4\p@ \@plus 1\p@   \@minus 2\p@}
\setlength{\partopsep    }{1\p@ \@plus 0.5\p@ \@minus 0.5\p@}
\setlength{\itemsep      }{2\p@ \@plus 1\p@   \@minus 0.5\p@}
\setlength{\parsep       }{2\p@ \@plus 1\p@   \@minus 0.5\p@}
\setlength{\leftmargin   }{3pc}
\setlength{\leftmargini  }{\leftmargin}
\setlength{\leftmarginii }{2em}
\setlength{\leftmarginiii}{1.5em}
\setlength{\leftmarginiv }{1.0em}
\setlength{\leftmarginv  }{0.5em}
\def\@listi  {\leftmargin\leftmargini}
\def\@listii {\leftmargin\leftmarginii
              \labelwidth\leftmarginii
              \advance\labelwidth-\labelsep
              \topsep  2\p@ \@plus 1\p@    \@minus 0.5\p@
              \parsep  1\p@ \@plus 0.5\p@ \@minus 0.5\p@
              \itemsep \parsep}
\def\@listiii{\leftmargin\leftmarginiii
              \labelwidth\leftmarginiii
              \advance\labelwidth-\labelsep
              \topsep    1\p@ \@plus 0.5\p@ \@minus 0.5\p@
              \parsep    \z@
              \partopsep 0.5\p@ \@plus 0\p@ \@minus 0.5\p@
              \itemsep \topsep}
\def\@listiv {\leftmargin\leftmarginiv
              \labelwidth\leftmarginiv
              \advance\labelwidth-\labelsep}
\def\@listv  {\leftmargin\leftmarginv
              \labelwidth\leftmarginv
              \advance\labelwidth-\labelsep}
\def\@listvi {\leftmargin\leftmarginvi
              \labelwidth\leftmarginvi
              \advance\labelwidth-\labelsep}

% create title
\providecommand{\maketitle}{}
\renewcommand{\maketitle}{%
  \par
  \begingroup
    \renewcommand{\thefootnote}{\fnsymbol{footnote}}
    % for perfect author name centering
    \renewcommand{\@makefnmark}{\hbox to \z@{$^{\@thefnmark}$\hss}}
    % The footnote-mark was overlapping the footnote-text,
    % added the following to fix this problem               (MK)
    \long\def\@makefntext##1{%
      \parindent 1em\noindent
      \hbox to 1.8em{\hss $\m@th ^{\@thefnmark}$}##1
    }
    \thispagestyle{empty}
    \@maketitle
    \@thanks
    \@notice
  \endgroup
  \let\maketitle\relax
  \let\thanks\relax
}

% rules for title box at top of first page
\newcommand{\@toptitlebar}{
  \hrule height 4\p@
  \vskip 0.25in
  \vskip -\parskip%
}
\newcommand{\@bottomtitlebar}{
  \vskip 0.29in
  \vskip -\parskip
  \hrule height 1\p@
  \vskip 0.09in%
}

% create title (includes both anonymized and non-anonymized versions)
\providecommand{\@maketitle}{}
\renewcommand{\@maketitle}{%
  \vbox{%
    \hsize\textwidth
    \linewidth\hsize
    \vskip 0.1in
    \@toptitlebar
    \centering
    {\LARGE\bf \@title\par}
    \@bottomtitlebar
    \if@anonymous
      \begin{tabular}[t]{c}\bf\rule{\z@}{24\p@}
        Anonymous Author(s) \\
        Affiliation \\
        Address \\
        \texttt{email} \\
      \end{tabular}%
    \else
      \def\And{%
        \end{tabular}\hfil\linebreak[0]\hfil%
        \begin{tabular}[t]{c}\bf\rule{\z@}{24\p@}\ignorespaces%
      }
      \def\AND{%
        \end{tabular}\hfil\linebreak[4]\hfil%
        \begin{tabular}[t]{c}\bf\rule{\z@}{24\p@}\ignorespaces%
      }
      \begin{tabular}[t]{c}\bf\rule{\z@}{24\p@}\@author\end{tabular}%
    \fi
    \vskip 0.3in \@minus 0.1in
  }
}

% add conference notice to bottom of first page
\newcommand{\ftype@noticebox}{8}
\newcommand{\@notice}{%
  % give a bit of extra room back to authors on first page
  \enlargethispage{2\baselineskip}%
  \@float{noticebox}[b]%
    \footnotesize\@noticestring%
  \end@float%
}

% abstract styling
\renewenvironment{abstract}%
{%
  \vskip 0.075in%
  \centerline%
  {\large\bf Abstract}%
  \vspace{0.5ex}%
  \begin{quote}%
}
{
  \par%
  \end{quote}%
  \vskip 1ex%
}

% For the paper checklist
\newcommand{\answerYes}[1][]{\textcolor{blue}{[Yes] #1}}
\newcommand{\answerNo}[1][]{\textcolor{orange}{[No] #1}}
\newcommand{\answerNA}[1][]{\textcolor{gray}{[NA] #1}}
\newcommand{\answerTODO}[1][]{\textcolor{red}{\bf [TODO]}}
\newcommand{\justificationTODO}[1][]{\textcolor{red}{\bf [TODO]}}

% handle tweaks for camera-ready copy vs. submission copy
\if@preprint
  \newcommand{\@noticestring}{%
    Preprint.%
  }
\else
  \if@neuripsfinal
    \newcommand{\@noticestring}{
      \@trackname
    }
  \else
    \newcommand{\@noticestring}{%
      Submitted to \@neuripsordinal\ Conference on Neural Information Processing Systems (NeurIPS \@neuripsyear). Do not distribute.%
    }

    % hide the acknowledgements
    \NewEnviron{hide}{}
    \let\ack\hide
    \let\endack\endhide

    % line numbers for submission
    \RequirePackage{lineno}
    \linenumbers

    % fix incompatibilities between lineno and amsmath, if required, by
    % transparently wrapping linenomath environments around amsmath
    % environments
    \AtBeginDocument{%
      \@ifpackageloaded{amsmath}{%
        \newcommand*\patchAmsMathEnvironmentForLineno[1]{%
          \expandafter\let\csname old#1\expandafter\endcsname\csname #1\endcsname
          \expandafter\let\csname oldend#1\expandafter\endcsname\csname end#1\endcsname
          \renewenvironment{#1}%
                          {\linenomath\csname old#1\endcsname}%
                          {\csname oldend#1\endcsname\endlinenomath}%
        }%
        \newcommand*\patchBothAmsMathEnvironmentsForLineno[1]{%
          \patchAmsMathEnvironmentForLineno{#1}%
          \patchAmsMathEnvironmentForLineno{#1*}%
        }%
        \patchBothAmsMathEnvironmentsForLineno{equation}%
        \patchBothAmsMathEnvironmentsForLineno{align}%
        \patchBothAmsMathEnvironmentsForLineno{flalign}%
        \patchBothAmsMathEnvironmentsForLineno{alignat}%
        \patchBothAmsMathEnvironmentsForLineno{gather}%
        \patchBothAmsMathEnvironmentsForLineno{multline}%
      }
      {}
    }
  \fi
\fi


\endinput
\end{filecontents*}

\usepackage[nonatbib]{neurips_2025}
\usepackage[utf8]{inputenc}
\usepackage[T1]{fontenc}
\usepackage{url}
\usepackage{microtype}
\usepackage{xcolor}

\usepackage{amsmath, amssymb, amsfonts}
\usepackage{booktabs}
\usepackage{multirow}
\usepackage{graphicx}
\usepackage{float}
\usepackage{placeins}
\usepackage{hyperref}

\hypersetup{
  colorlinks=true,
  linkcolor=blue,
  citecolor=blue,
  urlcolor=blue,
  pdftitle={RoPE Scaling as a Variational Inverse Problem: Exact Frequency Allocation and the Waterbed Trade-off},
  pdfauthor={Anonymous Authors}
}

\graphicspath{{./}}

% Robust figure inclusion: compile even when some legacy figure files are absent.
\newcommand{\safeincludegraphics}[2][]{%
  \IfFileExists{#2}{%
    \includegraphics[#1]{#2}%
  }{%
    \fbox{%
      \parbox[c][0.24\textheight][c]{0.92\linewidth}{%
        \centering Missing figure file: \texttt{\detokenize{#2}}%
      }%
    }%
  }%
}

% Math helpers
\newcommand{\Ediag}{E_{\mathrm{diag}}}
\newcommand{\Db}{\widehat{D}_b}
\newcommand{\sinc}{\mathrm{sinc}}
\DeclareMathOperator{\Ci}{Ci}

\newcommand{\asinh}{\mathrm{arcsinh}}

\title{RoPE Scaling as a Variational Inverse Problem: \\
Exact Frequency Allocation and the Waterbed Trade-off}

\author{
Anonymous Author(s) \\
Affiliation \\
Address \\
\texttt{email}
}
\date{}

\begin{document}
\maketitle

\begin{abstract}
RoPE context extension methods typically apply heuristic frequency warps without a principled objective.
We formulate RoPE frequency allocation as a joint variational problem that simultaneously minimizes broadband phase-collision energy and maximizes local Fisher resolution.
In the broadband limit, the full interference kernel decomposes into a diagonal ridge and a Brownian covariance, yielding a non-homogeneous governing ODE whose exact solution comprises a macroscopic hyperbolic tether ($\cosh$) and a microscopic Fisher resolution pulse ($b^{-2\phi}$).
This solution admits closed-form \emph{exact variational quantization} (EVQ) with a single physical parameter $\tau$ controlling the interference--resolution trade-off, and degrades exactly to standard geometric RoPE as $\tau\to 0$.
Empirically, from-scratch TinyStories scaling (50M--350M) validates the predicted high-frequency bias.
Controlled-protocol LoRA suites on Llama-3-8B and Qwen-2.5-7B with dual-seed validation ($42$/$1337$) reveal a structured waterbed trade-off: EVQ scheduling improves retrieval-oriented tasks while regressing on multi-hop reasoning, with both directions reproduced across seeds and model families---an outcome quantitatively predicted by the waterbed inequality.
We characterize \emph{when and why} spectral reshaping helps or hurts, providing a principled framework and actionable regime map for RoPE scaling.
\end{abstract}

\section{Introduction}
Rotary positional embedding (RoPE)~\cite{su2021roformer} parameterizes relative positions through a bank of sinusoidal frequencies, enabling efficient extrapolation beyond the training context. Because long-context capability is now a bottleneck for retrieval-augmented generation, tool use, and code agents, a growing ecosystem of RoPE-based context extension methods has emerged: position interpolation (PI)~\cite{chen2023positioninterpolation}, YaRN~\cite{peng2023yarn}, base-adjustment heuristics, and warp variants such as LongRoPE~\cite{ding2024longrope}. Yet it remains unclear what these frequency manipulations are optimizing and why particular shapes work across models and tasks.

This paper offers a theory-first answer. We formulate RoPE frequency allocation as a \emph{joint} variational problem: minimize phase-collision energy (long-range interference) while maximizing local Fisher resolution (short-range discriminability), subject to a task-dependent distance prior $D(\Delta)$. In the broadband limit, the interference kernel separates into a diagonal ridge and a Brownian covariance. The resulting Euler--Lagrange equation is a non-homogeneous ODE whose exact solution reveals two competing forces: a macroscopic hyperbolic tether for interference suppression and a microscopic exponential pulse for local resolution. This duality provides a precise mathematical formulation of the ``waterbed'' trade-off inherent in all RoPE scaling.

\textbf{Contributions.}
\begin{itemize}
  \item \textbf{Joint variational framework with governing ODE:} we derive a non-homogeneous ODE for RoPE frequency density from a joint interference--Fisher objective. Its exact solution (Theorem~1) decomposes into a $\cosh$-type broadband tether and a $b^{-2\phi}$ Fisher pulse, with a single physical parameter $\tau$ governing the trade-off.
  \item \textbf{Exact Variational Quantization (EVQ):} we obtain a closed-form analytical warp $\phi_k(\tau)=1-\tau^{-1}\asinh\bigl((1-u_k)\sinh\tau\bigr)$ that maps continuous density to discrete frequencies with zero heuristic parameters. This formula degrades exactly to geometric RoPE as $\tau\to 0$ (Theorem~2), unifying pre-training and context extension.
  \item \textbf{Cross-scale, cross-model waterbed verification:} from-scratch TinyStories scaling (50M--350M) validates predicted high-frequency bias; controlled-protocol dual-seed LoRA evaluation on Llama-3-8B and Qwen-2.5-7B with shared \texttt{inv\_freq.copy()} injection reveals a structured task-family trade-off---retrieval gains coexist with multi-hop losses---quantitatively consistent with the waterbed inequality.
\end{itemize}

\section{Related Work}
\paragraph{Rotary and relative position encodings.}
RoPE~\cite{su2021roformer} implements relative position via complex rotations; other relative encodings include ALiBi~\cite{press2021alibi}, Transformer-XL~\cite{dai2019transformerxl}, and T5-style relative bias~\cite{raffel2020t5}. Alternative constructions such as Kerple~\cite{chi2022kerple} and FIRE~\cite{li2023fire} target extrapolation or inductive biases.

\paragraph{Context window extension.}
PI~\cite{chen2023positioninterpolation} compresses positions; YaRN~\cite{peng2023yarn} uses progressive warping; community NTK-aware scaling heuristics modify the base to better preserve kernel behavior; and later variants such as LongRoPE~\cite{ding2024longrope} explore additional warps. Orthogonal approaches extend context via sparse attention: Sparse Transformers~\cite{child2019sparse}, Longformer~\cite{beltagy2020longformer}, BigBird~\cite{zaheer2020bigbird}, and Reformer~\cite{kitaev2020reformer}.

\paragraph{Theoretical analyses.}
Recent work studies how positional design affects length generalization~\cite{kazemnejad2023impact}. Our analysis complements this line by treating scaling as a variational inverse problem with explicit priors and structural conditions.

\paragraph{Scaling and evaluation.}
We follow standard scaling-law and evaluation practices~\cite{kaplan2020scaling,hoffmann2022computeoptimal}. For long-context evaluation we use LongBench~\cite{bai2023longbench} and NIAH~\cite{kamradt2023niah}. Our 8B suite uses LoRA~\cite{hu2022lora} and efficient attention kernels~\cite{dao2022flashattention,dao2023flashattention2}.

\paragraph{How this work differs.}
This work provides a unified structural lens: given a distance prior, what RoPE frequency density is optimal, and when can proxies mislead?

\section{A Joint Variational Framework for RoPE Scaling}

\subsection{RoPE as frequency allocation}
Let $d$ be the head dimension and $N=d/2$ the number of RoPE frequencies. Standard geometric RoPE uses $\omega_k=b^{-k/(N-1)}$ for base $b$. We pass to a continuous coordinate $\phi\in[0,1]$ so that $\omega(\phi)=b^{-\phi}$ and a nonnegative density $\rho(\phi)$ encodes how dimensions are allocated along $\phi$ with $\int_0^1 \rho(\phi)\,d\phi = 1$. A relative distance $\Delta$ induces a phase similarity $S_{\rho}(\Delta) = \int_0^1 \rho(\phi)\cos(\omega(\phi)\,\Delta)\,d\phi$.

\subsection{Phase-collision energy and broadband scale separation}
The quadratic phase-collision energy under a distance prior $D(\Delta)$ is
\begin{equation}
\mathcal{C}[\rho] = \int D(\Delta)\,S_{\rho}(\Delta)^2\,d\Delta = \iint \rho(\phi_1)\rho(\phi_2)\,K(\phi_1,\phi_2)\,d\phi_1\,d\phi_2,
\label{eq:full_C}
\end{equation}
where $K(\phi_1,\phi_2)=\int D(\Delta)\cos(\omega(\phi_1)\Delta)\cos(\omega(\phi_2)\Delta)\,d\Delta$ is the induced kernel.

In the broadband limit ($b\to\infty$, power-law tail prior $D(\Delta)\propto 1/(\Delta\ln L)$), the kernel undergoes a strict \emph{scale separation}: the self-interference ($\phi_1=\phi_2$) produces a singular diagonal ridge, while the cross-frequency coupling ($\phi_1\neq\phi_2$) converges to a Brownian covariance (Appendix~F):
\begin{equation}
K(\phi_1,\phi_2) \;\xrightarrow{b,L\to\infty}\; \alpha\,\delta(\phi_1-\phi_2) + \beta\,\min(\phi_1,\phi_2),
\label{eq:scale_sep}
\end{equation}
where $\alpha>0$ controls the local diagonal ridge strength and $\beta>0$ the global cross-frequency coupling. Previous diagonal approximations~\cite{su2021roformer} retained only the $\alpha$-term, discarding the $\beta$-term that typically accounts for $\sim\!11\%$ of the energy at finite base. The decomposition \eqref{eq:scale_sep} captures both components, yielding the interference functional
\begin{equation}
\mathcal{C}_{\mathrm{interf}}[\rho] = \frac{\alpha}{2}\int_0^1 \rho(\phi)^2\,d\phi + \frac{\beta}{2}\iint_0^1 \rho(\phi_1)\rho(\phi_2)\min(\phi_1,\phi_2)\,d\phi_1\,d\phi_2.
\label{eq:C_interf}
\end{equation}

\subsection{Joint objective with Fisher fidelity}
Minimizing $\mathcal{C}_{\mathrm{interf}}$ alone would collapse density toward high frequencies, destroying local position resolution. For a RoPE encoding, the Fisher information for estimating relative position $\Delta$ from frequency coordinate $\phi$ scales as $\omega(\phi)^2=b^{-2\phi}$; the aggregate local resolution is $\mathcal{I}[\rho]=\int_0^1 \rho(\phi)\,b^{-2\phi}\,d\phi$. We therefore construct the \emph{joint variational objective}:
\begin{equation}
\mathcal{J}[\rho]
= \mathcal{C}_{\mathrm{interf}}[\rho]
- \mu\!\int_0^1\!\rho(\phi)\,b^{-2\phi}\,d\phi
+ \lambda\!\left(\int_0^1\!\rho(\phi)\,d\phi - 1\right),
\label{eq:joint_obj}
\end{equation}
where $\mu\ge 0$ weights the Fisher fidelity term and $\lambda$ enforces normalization.

Setting $\delta\mathcal{J}/\delta\rho=0$ gives the stationarity condition
\begin{equation}
\alpha\rho(\phi)+\beta\!\int_0^1\!\rho(\psi)\min(\phi,\psi)\,d\psi - \mu\,b^{-2\phi}+\lambda=0.
\label{eq:stationarity}
\end{equation}
Differentiating once eliminates $\lambda$ and produces $\int_\phi^1\rho(\psi)\,d\psi$; differentiating again eliminates the integral. Defining $\tau^2=\beta/\alpha$ and $\gamma_F=4\mu(\ln b)^2/\alpha$, we obtain the \textbf{governing non-homogeneous ODE}:
\begin{equation}
\rho''(\phi)-\tau^2\rho(\phi)=\gamma_F\,b^{-2\phi}.
\label{eq:governing_ode}
\end{equation}
The full derivation, including boundary conditions from \eqref{eq:stationarity}, is in Appendix~A.

\textbf{Theorem 1 (Exact solution of the joint variational ODE).}
\textit{The general solution of \eqref{eq:governing_ode} is}
\begin{equation}
\rho^{\star}(\phi) = \underbrace{C_1\cosh(\tau\phi)+C_2\sinh(\tau\phi)}_{\text{broadband hyperbolic tether}} + \underbrace{\frac{\gamma_F}{4(\ln b)^2-\tau^2}\;b^{-2\phi}}_{\text{Fisher resolution pulse}},
\label{eq:joint_solution}
\end{equation}
\textit{where $C_1,C_2$ are fixed by the boundary condition $\alpha\rho'(1)+2\mu(\ln b)b^{-2}=0$ and normalization $\int_0^1\rho\,d\phi=1$, provided $\tau\neq 2\ln b$.}

\textbf{Proof sketch.} The homogeneous equation $\rho''-\tau^2\rho=0$ has basis $\{\cosh(\tau\phi),\sinh(\tau\phi)\}$. The particular solution is obtained by substituting $\rho_p=Ab^{-2\phi}$, giving $A(4(\ln b)^2-\tau^2)=\gamma_F$. Boundary conditions follow from evaluating \eqref{eq:stationarity} and its derivative at $\phi=1$. \hfill$\square$

\textbf{Physical interpretation.} The solution \eqref{eq:joint_solution} decomposes optimal frequency allocation into two competing forces: (i) the \emph{hyperbolic tether} $C_1\cosh(\tau\phi)+C_2\sinh(\tau\phi)$, shaped by global cross-frequency coupling, which distributes density to suppress long-range phase collisions; and (ii) the \emph{Fisher pulse} $\propto b^{-2\phi}$, which concentrates density at high frequencies ($\phi\approx 0$) to preserve local position resolution. The balance between these forces is governed by a single physical parameter $\tau=\sqrt{\beta/\alpha}$, the ratio of global coupling to local ridge strength.

When $\mu=0$ (pure interference minimization), the Fisher pulse vanishes and the solution reduces to $\rho\propto\cosh(\tau(1-\phi))$ under the natural boundary condition $\rho'(1)=0$---recovering the broadband result of the pure-interference analysis (Appendix~F).

\subsection{Exact Variational Quantization (EVQ)}
\label{sec:evq}
Existing warp heuristics (piecewise-linear in YaRN, sigmoid curves, etc.) approximate the optimal density through multi-parameter curve fitting. These introduce artificial inflection points that violate the strict convexity ($\rho''>0$) of the theoretical optimum, amplifying discretization error (Appendix~G, Proposition~G.5).

We instead perform \emph{exact variational quantization}: compute the CDF of the dominant $\cosh$ component analytically and invert it in closed form. The CDF of $\rho(\phi)\propto\cosh(\tau(1-\phi))$ is
\begin{equation}
F(\phi)=1-\frac{\sinh(\tau(1-\phi))}{\sinh(\tau)}.
\label{eq:evq_cdf}
\end{equation}
Setting $F(\phi_k)=u_k=(k+\tfrac{1}{2})/N$ and inverting yields the \textbf{closed-form EVQ warp}:
\begin{equation}
\boxed{\;\phi_k(\tau) = 1-\frac{1}{\tau}\,\asinh\!\Big((1-u_k)\sinh(\tau)\Big),\qquad \omega_k=b^{-\phi_k(\tau)},\quad k=0,\dots,N{-}1.\;}
\label{eq:evq_warp}
\end{equation}
This formula has a \emph{single} physical parameter $\tau$, the prior interference coupling ratio. No curve-fitting, grid search, or heuristic shape parameters are involved.

\textbf{Theorem 2 (Asymptotic degradation to geometric RoPE).}
\textit{As $\tau\to 0^+$, the EVQ schedule degrades smoothly to the uniform grid:}
\begin{equation}
\lim_{\tau\to 0}\phi_k(\tau)=u_k=\frac{k+\tfrac{1}{2}}{N},
\label{eq:degradation}
\end{equation}
\textit{corresponding to standard geometric RoPE.}

\textbf{Proof.} For small $\tau$, $\sinh(\tau)\approx\tau$ and $\asinh(x\tau)\approx x\tau$. Then $\phi_k\approx 1-(1-u_k)=u_k$. \hfill$\square$

This theorem establishes that geometric RoPE is not an alternative to the variational framework but its $\tau=0$ special case. Context extension corresponds to increasing $\tau>0$ to redistribute density toward high frequencies, with the trade-off governed by the waterbed inequality (Section~4).

\subsection{Structural theorems for the diagonal regime}
The diagonal component of $K$ (the $\alpha$-term in \eqref{eq:scale_sep}) admits $\Ediag(\phi)=\frac{1}{2}(1+\Db(2b^{-\phi}))$ where $\Db(\xi)=\int D(\Delta)\cos(\xi\Delta)\,d\Delta$, with inverse-law optimizer $\rho^{\star}\propto 1/\Ediag(\phi)$. Three structural results characterize its prior-dependence (proofs in Appendices~B--D):

\textbf{Proposition 1 (Uniform prior $\Rightarrow$ geometric).}
\textit{For $D(\Delta)=\frac{1}{L}\mathbb{1}_{[0,L]}$ with $L\to\infty$: $\Ediag\to\tfrac{1}{2}$ uniformly, so $\rho^{\star}\to 1$ (geometric).}

\textbf{Proposition 2 (Power-law $\Rightarrow$ high-frequency bias).}
\textit{For $D(\Delta)\propto 1/(\Delta\ln L)$ on $[1,L]$: $\Ediag(\phi)\approx \frac{1}{2}(1+(\phi\ln b-\gamma-\ln 2)/\ln L)$ is increasing, so $\rho^{\star}$ is decreasing (high-frequency biased) with bounded amplitude ratio $1+\mathcal{O}(\ln b/\ln L)$.}

\textbf{Proposition 3 (Bimodal prior $\Rightarrow$ resonance trap).}
\textit{For $D=\lambda\delta(\Delta-s)+(1-\lambda)\delta(\Delta-\ell)$: near-resonant frequency alignments create $\mathcal{O}(\varepsilon^2)$-deep diagonal minima, making the inverse-law solution brittle.}

Proposition~1 connects to Theorem~2: both predict geometric RoPE under task-agnostic or local-only priors. Proposition~3 motivates the bounded Fisher pulse in Theorem~1 as a regularizer against resonance-driven density collapse.

\section{Predictions}

\subsection{Waterbed inequality}
Define local Fisher information $\mathcal{I}(\phi)=c\,\rho(\phi)\,b^{-2\phi}$ and local error proxy $E(\phi)$. A Jensen-based argument yields the information-theoretic scaling bound
\begin{equation}
\int_0^1 \ln E(\phi)\,d\phi \ge \ln b-\ln c,
\label{eq:waterbed}
\end{equation}
proved in Appendix~F. This is an integrated log-error lower bound: increasing base $b$ necessarily incurs error somewhere, formalizing a strict zero-sum trade-off between phase-collision suppression and local Fisher resolution. This inequality is the mathematical backbone of the task-family trade-offs observed empirically in Section~5.

\subsection{Finite-base calibration and \texorpdfstring{$L/b$}{L/b} transition}
Two numerical checks tie the theory to finite-base regimes. First, for $L=16384$ and $b\in\{10^4,5\times10^5\}$, $\Ediag(\phi)$ is accurately captured by the Proposition~2 affine form ($R^2_{\text{mid}}\approx 0.994,0.995$; Table~\ref{tab:affine_calibration}). Second, a scan over prior mixtures reveals a regime transition: at low $L/b\in\{1.6,10\}$, EVQ keeps lower phase-collision energy than geometric across the scan; at high $L/b\in\{100,1000\}$, a low-$p$ crossing appears. This operationalizes the regime-dependent nature of spectral reshaping.

\begin{table}[t]
\centering
\begin{tabular}{lcccc}
\toprule
Case & $B_{\text{fit}}$ & $B_{\text{th}}$ & $R^2_{\text{mid}}$ & max err. \\
\midrule
$b=10^4,\;L=16384$ & 0.4746 & 0.4746 & 0.9942 & 0.0759 \\
$b=5\times10^5,\;L=16384$ & 0.6724 & 0.6761 & 0.9954 & 0.0428 \\
\bottomrule
\end{tabular}
\caption{Finite-base affine calibration. Close slope match and high mid-band $R^2$ support the affine approximation in the tested range.}
\label{tab:affine_calibration}
\end{table}

\subsection{Implicit priors of existing methods}
PI~\cite{chen2023positioninterpolation} applies uniform compression $\Delta\to\Delta/\alpha$, corresponding to $\rho_{\mathrm{PI}}=1$ with rescaled base---implicitly assuming a self-similar distance prior. YaRN~\cite{peng2023yarn} implements a piecewise warp, corresponding to a step-function density with extra low-frequency mass. Our framework unifies these as special cases: each method implicitly selects a density $\rho$ and, by the inverse law, an implicit distance prior $D_{\mathrm{implicit}}$. EVQ makes this selection explicit through $\tau$: $\tau=0$ recovers geometric (Theorem~2), and increasing $\tau$ progressively tilts density toward high frequencies at a rate governed by the waterbed inequality. Proposition~3 warns that sharp band concentration can produce brittle resonance traps, motivating the bounded Fisher pulse in Theorem~1 as a natural regularizer.

\section{Experiments}

\subsection{Protocol summary}
We report (i)~from-scratch TinyStories training~\cite{eldan2023tinystories} with controlled steps/tokens and PPL measured at multiple lengths, and (ii)~controlled-protocol LoRA studies on Llama-3-8B and Qwen-2.5-7B: identical budgets, a single injection path (\texttt{inv\_freq.copy()}), and no built-in \texttt{rope\_scaling} API differences across methods (Appendix~H). The EVQ schedule~\eqref{eq:evq_warp} is implemented with $\tau$ chosen per model family; for TinyStories pre-training, an equivalent density-guided schedule is used. Unless noted, results are single-seed for large models and multi-seed for 50M.

\subsection{From-scratch TinyStories scaling (50M--350M)}

\begin{table}[t]
\centering
\begin{tabular}{lccc}
\toprule
Model & Geometric (2K / 16K) & EVQ (2K / 16K) & $\Delta@16\mathrm{K}$ \\
\midrule
50M (3 seeds) & $6.821\pm0.067$ / $19.386\pm2.007$ & $6.686\pm0.200$ / $17.404\pm1.564$ & $-10.2\%$ \\
100M & $6.457$ / $10.888$ & $6.382$ / $9.417$ & $-13.5\%$ \\
350M & $2.477\pm0.131$ / $14.653\pm3.851$ & $2.467\pm0.126$ / $12.646\pm3.093$ & $-13.7\%$ \\
\bottomrule
\end{tabular}
\caption{From-scratch scaling on TinyStories. PPL at train length (2K) and extrapolation length (16K). EVQ uses the theory-derived high-frequency-biased schedule. Reproducibility settings in Appendix~H.}
\label{tab:tinystories_scaling}
\end{table}

Table~\ref{tab:tinystories_scaling} shows that the EVQ-derived density improves long-context perplexity consistently across scales ($-10\%$ to $-14\%$ at 16K) while preserving short-context fit, validating the predicted high-frequency bias of Proposition~2. A supplementary YaRN comparison confirms that YaRN's implicit prior is mismatched to TinyStories extrapolation (Appendix~I).

\subsection{Controlled-protocol LoRA evaluation}

\textbf{Llama-3-8B (5 schedules, 600 steps).}
All methods use minimally adapted base models with LoRA on WikiText-style corpora, without instruction alignment. LongBench is reported on its native $[0,1]$ scale; relative differences are the signal of interest.

\begin{table}[t]
\centering
\begin{tabular}{lccc}
\toprule
Method & NIAH@16K & LongBench mean & Passkey@16K \\
\midrule
Geometric ($\tau{=}0$) & 0.9545 & 0.0626 & 4.6846 \\
PI & 1.0000 & 0.0665 & 6.5850 \\
YaRN & 1.0000 & 0.0656 & 5.2433 \\
Diagonal-only & 1.0000 & 0.0687 & 5.4137 \\
EVQ ($\tau{>}0$) & 1.0000 & 0.0717 & 6.5242 \\
\bottomrule
\end{tabular}
\caption{Llama-3-8B controlled protocol evaluation ($n{=}6$ LongBench tasks). EVQ achieves the highest LongBench mean ($+14.5\%$ vs geometric). ``Diagonal-only'' uses the inverse-law solution without Fisher regularization; EVQ adds the bounded Fisher pulse. All methods share the same \texttt{inv\_freq.copy()} path. Appendix~H provides task-level provenance.}
\label{tab:8b_fair}
\end{table}

EVQ achieves the strongest LongBench mean among five schedules. Compared with the diagonal-only variant, EVQ improves LongBench by $+4.4\%$ and Passkey by $\sim\!20\%$, directly validating that the Fisher pulse in Theorem~1 contributes measurable gains beyond pure interference minimization. Task-level paired analysis ($n=6$) yields non-significant $p$-values (Appendix~H), consistent with mechanism validation rather than leaderboard claims.

\textbf{Qwen-2.5-7B (dual seed, 400 steps, LongBench-21).}
To test cross-model generality, we replicate on Qwen-2.5-7B-Instruct with seeds $42$ and $1337$ (Table~\ref{tab:qwen_dual_seed}).

\begin{table}[t]
\centering
\begin{tabular}{lcccc}
\toprule
Seed & Geometric & EVQ & $\Delta$ & Record (W/T/L) \\
\midrule
42 & 44.44 & 44.08 & $-0.35$ & 8 / 2 / 11 \\
1337 & 44.47 & 44.05 & $-0.42$ & 9 / 0 / 12 \\
\midrule
Mean & 44.45 & 44.07 & $-0.39$ & --- \\
\bottomrule
\end{tabular}
\caption{Qwen-2.5-7B LongBench-21 under dual-seed protocol. Aggregate delta is negative ($p_{\mathrm{FDR}}=0.010$), but task-family decomposition reveals structured waterbed behavior.}
\label{tab:qwen_dual_seed}
\end{table}

\subsection{Empirical verification of the waterbed trade-off}
\label{sec:waterbed_empirical}
The aggregate Qwen score shows a small overall regression ($-0.39$ points, $p_{\mathrm{FDR}}=0.010$). Decomposing by task family (Table~\ref{tab:qwen_task_family}) reveals a \emph{consistent structured trade-off} reproduced across both seeds:

\begin{table}[t]
\centering
\begin{tabular}{lcc}
\toprule
Task Family & Representative Tasks & Dual-Seed Avg.\ $\Delta$ \\
\midrule
Retrieval & passage\_retrieval\_en & $+2.50$ \\
Single-doc QA & NarrativeQA, TriviaQA & $+0.84$ \\
Summarization & QMSum, MultiNews & $\approx 0$ \\
Multi-hop QA & MuSiQue, 2WikiMQA, HotpotQA & $-2.69$ \\
Code & RepoBench-P & $-1.42$ \\
\bottomrule
\end{tabular}
\caption{Task-family decomposition, Qwen-2.5-7B EVQ vs.\ geometric (dual-seed average). The retrieval--gain / multi-hop--loss pattern reproduces across seeds and matches the Llama-3-8B directional structure.}
\label{tab:qwen_task_family}
\end{table}

\textbf{Waterbed interpretation.} The observed polarization---retrieval gains coexisting with multi-hop losses---provides direct empirical verification of the waterbed inequality \eqref{eq:waterbed} and the dual-force structure of Theorem~1:

\emph{(i) Fisher pulse decay at extreme base.} In the exact solution \eqref{eq:joint_solution}, the Fisher resolution pulse decays as $b^{-2\phi}$. Qwen-2.5 uses base $b=10^6$, so the pulse decays as $10^{-12\phi}$---orders of magnitude faster than at the conventional $b=10^4$ used in Llama-3. Under this extreme decay rate, redistributing mid-frequency density to suppress long-range collisions rapidly depletes the Fisher information available for local token-level reasoning, precisely the resource required by multi-hop tasks such as MuSiQue ($\Delta=-3.25$).

\emph{(ii) Prior mismatch.} The 400-step WikiText LoRA provides unimodal local gradient signal. Under the EVQ-reshaped frequency manifold, deeper attention heads receive no multi-hop cross-frequency supervision, resulting in manifold misalignment on bimodal-prior tasks (cf.\ Proposition~3).

\emph{(iii) Cross-model consistency.} The Llama-3-8B suite shows the same directional pattern: 4/6 task wins for EVQ, with losses on multi-step reasoning tasks. The Qwen validation ($n=21$ tasks, dual seed) confirms this is a structural consequence of spectral reshaping, not an implementation artifact.

These observations support the central thesis: context extension is a strict zero-sum trade-off between long-range interference suppression and local Fisher resolution, governed by the single parameter $\tau$ and bounded by the waterbed inequality. Task-agnostic ``universal improvement'' is a mathematical impossibility when the true distance prior $D(\Delta)$ is multi-modal.

\FloatBarrier

\section{Limitations}
\begin{itemize}
  \item \textbf{Scale-separation approximation.} The broadband decomposition $K\to\alpha\delta+\beta\min$ is asymptotically exact as $b\to\infty$ but introduces $\mathcal{O}(1/\ln b)$ residuals at finite base (Appendix~E). For small $b$ the governing ODE may underestimate cross-frequency coupling.

  \item \textbf{Waterbed trade-off is intrinsic, not defeatable.} At the 7B/8B scale, EVQ does not uniformly dominate baselines: it improves retrieval-oriented tasks but regresses on multi-hop reasoning (Tables~\ref{tab:qwen_dual_seed}--\ref{tab:qwen_task_family}). The aggregate Qwen delta is $-0.39$ points with $p_{\mathrm{FDR}}=0.010$. We report this as a structural consequence of the waterbed inequality rather than a limitation to overcome.

  \item \textbf{Statistical power at task level.} The Llama-3-8B LongBench suite ($n=6$ tasks) has limited statistical power; the Qwen-2.5-7B expansion to $n=21$ tasks provides stronger task-family decomposition but relies on a single controlled protocol (WikiText LoRA, 400 steps). Different fine-tuning regimes may shift the trade-off boundary.

  \item \textbf{Coverage boundaries.} The current protocol covers geometric/PI/YaRN/diagonal-only/EVQ on Llama-3-8B and EVQ vs.\ baseline on Qwen-2.5-7B; a matched LongRoPE run and additional model families are pending.
\end{itemize}

\section{Conclusion}
We have recast RoPE context extension as a variational inverse problem. A joint phase-collision--Fisher-fidelity functional, after broadband scale separation, yields a non-homogeneous governing ODE whose exact solution (Theorem~1) is a cosh tether modulated by a Fisher pulse at rate $b^{-2\phi}$. The solution discretises into EVQ---a closed-form $\asinh$ warp parameterised by a single physical quantity $\tau=\sqrt{\beta/\alpha}$ that degrades exactly to geometric RoPE as $\tau\to 0$ (Theorem~2).

Empirically, from-scratch TinyStories scaling (50M--350M) validates the predicted high-frequency bias for pre-training. At the 7B/8B scale, controlled-protocol LoRA evaluation on two model families (Llama-3 with $b\approx 10^4$, Qwen-2.5 with $b\approx 10^6$) with dual-seed validation reveals a \emph{structured trade-off}: EVQ consistently improves retrieval-oriented tasks while incurring losses on multi-hop reasoning, with the pattern reproduced across seeds and models. This trade-off is a direct, quantitative consequence of the waterbed inequality; the stronger decay of the Fisher pulse at large $b$ correctly predicts the reduced EVQ margin on Qwen.

The practical takeaway is \emph{task-aware}: spectral reshaping is most beneficial when the downstream distribution is retrieval-heavy, while multi-hop reasoning workloads may prefer conservative schedules closer to geometric. Future work includes task-conditional priors that adapt $\tau$ to workload mixtures and layer-wise frequency allocation.
Broader-impact considerations are provided after References.

\clearpage
\phantomsection
\label{sec:references_start}
\begin{thebibliography}{34}

\bibitem{su2021roformer}
J.~Su, Y.~Lu, S.~Pan, A.~Wen, and Y.~Liu.
\newblock RoFormer: Enhanced Transformer with Rotary Position Embedding.
\newblock \emph{arXiv preprint arXiv:2104.09864}, 2021.

\bibitem{press2021alibi}
O.~Press, N.~A. Smith, and M.~Lewis.
\newblock Train Short, Test Long: Attention with Linear Biases Enables Input Length Extrapolation.
\newblock \emph{arXiv preprint arXiv:2108.12409}, 2021.

\bibitem{dai2019transformerxl}
Z.~Dai, Z.~Yang, Y.~Yang, J.~Carbonell, Q.~V. Le, and R.~Salakhutdinov.
\newblock Transformer-XL: Attentive Language Models Beyond a Fixed-Length Context.
\newblock In \emph{ACL}, 2019.

\bibitem{raffel2020t5}
C.~Raffel et~al.
\newblock Exploring the Limits of Transfer Learning with a Unified Text-to-Text Transformer.
\newblock \emph{JMLR}, 2020.

\bibitem{chi2022kerple}
E.~Chi, H.~Zhang, and others.
\newblock Kerple: Kernelized Relative Position Encoding for Length Extrapolation.
\newblock \emph{arXiv preprint arXiv:2205.09921}, 2022.

\bibitem{li2023fire}
Y.~Li et~al.
\newblock FIRE: Flexible In-context RoPE Enhancement for Length Generalization.
\newblock \emph{arXiv preprint arXiv:2308.02800}, 2023.

\bibitem{chen2023positioninterpolation}
C.~Chen, J.~Liu, and others.
\newblock Extending Context Window of Large Language Models via Position Interpolation.
\newblock \emph{arXiv preprint arXiv:2306.15595}, 2023.

\bibitem{peng2023yarn}
B.~Peng, H.~Zhang, and others.
\newblock YaRN: Efficient Context Window Extension of Large Language Models.
\newblock \emph{arXiv preprint arXiv:2309.00071}, 2023.

\bibitem{ding2024longrope}
J.~Ding, P.~Zhang, and others.
\newblock LongRoPE: Extending LLM Context Window Beyond 2 Million Tokens.
\newblock \emph{arXiv preprint arXiv:2402.13753}, 2024.

\bibitem{child2019sparse}
R.~Child, S.~Gray, A.~Radford, and I.~Sutskever.
\newblock Generating Long Sequences with Sparse Transformers.
\newblock \emph{arXiv preprint arXiv:1904.10509}, 2019.

\bibitem{beltagy2020longformer}
I.~Beltagy, M.~E. Peters, and A.~Cohan.
\newblock Longformer: The Long-Document Transformer.
\newblock \emph{arXiv preprint arXiv:2004.05150}, 2020.

\bibitem{zaheer2020bigbird}
M.~Zaheer et~al.
\newblock BigBird: Transformers for Longer Sequences.
\newblock In \emph{NeurIPS}, 2020.

\bibitem{kitaev2020reformer}
N.~Kitaev, L.~Kaiser, and A.~Levskaya.
\newblock Reformer: The Efficient Transformer.
\newblock In \emph{ICLR}, 2020.

\bibitem{kazemnejad2023impact}
A.~Kazemnejad et~al.
\newblock The Impact of Positional Encoding on Length Generalization in Transformers.
\newblock In \emph{NeurIPS}, 2023.

\bibitem{kaplan2020scaling}
J.~Kaplan et~al.
\newblock Scaling Laws for Neural Language Models.
\newblock \emph{arXiv preprint arXiv:2001.08361}, 2020.

\bibitem{hoffmann2022computeoptimal}
J.~Hoffmann et~al.
\newblock Training Compute-Optimal Large Language Models.
\newblock \emph{arXiv preprint arXiv:2203.15556}, 2022.

\bibitem{bai2023longbench}
Y.~Bai et~al.
\newblock LongBench: A Bilingual, Multitask Benchmark for Long Context Understanding.
\newblock \emph{arXiv preprint arXiv:2308.14508}, 2023.

\bibitem{kamradt2023niah}
G.~Kamradt.
\newblock Needle-in-a-Haystack: Evaluating Long-Context Retrieval (technical report), 2023.

\bibitem{hu2022lora}
E.~Hu, Y.~Shen, P.~Wallis, and others.
\newblock LoRA: Low-Rank Adaptation of Large Language Models.
\newblock In \emph{ICLR}, 2022.

\bibitem{dao2022flashattention}
T.~Dao, D.~Fu, S.~Ermon, A.~R{\'e}, and C.~Rudin.
\newblock FlashAttention: Fast and Memory-Efficient Exact Attention with IO-Awareness.
\newblock In \emph{NeurIPS}, 2022.

\bibitem{dao2023flashattention2}
T.~Dao.
\newblock FlashAttention-2: Faster Attention with Better Parallelism and Work Partitioning.
\newblock \emph{arXiv preprint arXiv:2307.08691}, 2023.

\bibitem{eldan2023tinystories}
R.~Eldan and Y.~Li.
\newblock TinyStories: How Small Can Language Models Be and Still Speak Coherent English?
\newblock \emph{arXiv preprint arXiv:2305.07759}, 2023.

\bibitem{grattafiori2024llama3}
A.~Grattafiori et~al.
\newblock The Llama 3 Herd of Models.
\newblock \emph{arXiv preprint arXiv:2407.21783}, 2024.

\bibitem{touvron2023llama}
H.~Touvron et~al.
\newblock LLaMA: Open and Efficient Foundation Language Models.
\newblock \emph{arXiv preprint arXiv:2302.13971}, 2023.

\bibitem{touvron2023llama2}
H.~Touvron et~al.
\newblock LLaMA 2: Open Foundation and Fine-Tuned Chat Models.
\newblock \emph{arXiv preprint arXiv:2307.09288}, 2023.

\bibitem{roziere2023codellama}
B.~Rozi{\`e}re et~al.
\newblock Code Llama: Open Foundation Models for Code.
\newblock \emph{arXiv preprint arXiv:2308.12950}, 2023.

\bibitem{vaswani2017attention}
A.~Vaswani et~al.
\newblock Attention Is All You Need.
\newblock In \emph{NeurIPS}, 2017.

\bibitem{cover2006elements}
T.~M. Cover and J.~A. Thomas.
\newblock \emph{Elements of Information Theory}.
\newblock Wiley, 2nd edition, 2006.

\bibitem{vantrees1968detection}
H.~L. Van Trees.
\newblock \emph{Detection, Estimation, and Modulation Theory}.
\newblock Wiley, 1968.

\end{thebibliography}

\section*{Broader Impact}
This work is a methodological framework for long-context positional design. Potential positive impact is improved stability diagnostics for context extension, which can reduce trial-and-error retraining cost when adapting models to longer windows. Potential negative impact is that longer-context capability can also strengthen misuse scenarios (e.g., scalable synthesis of harmful content and broader retrieval of sensitive information). Our contribution does not release new high-risk models or new datasets; nevertheless, we recommend standard deployment controls for long-context systems, including access restriction, content safeguards, and monitoring for privacy-sensitive prompts.

\section*{NeurIPS Paper Checklist}
\begin{enumerate}

\item {\bf Claims}
    \item[] Question: Do the main claims made in the abstract and introduction accurately reflect the paper's contributions and scope?
    \item[] Answer: \answerYes{}
    \item[] Justification: The abstract and introduction explicitly frame the contribution as a joint variational framework with exact ODE solution, EVQ closed-form warp, and waterbed trade-off verification. Claims are aligned with reported statistics and stated limitations.

\item {\bf Limitations}
    \item[] Question: Does the paper discuss the limitations of the work performed by the authors?
    \item[] Answer: \answerYes{}
    \item[] Justification: A dedicated Limitations section discusses scale-separation residuals, the intrinsic waterbed trade-off, statistical power boundaries, and baseline coverage gaps.

\item {\bf Theory assumptions and proofs}
    \item[] Question: For each theoretical result, does the paper provide the full set of assumptions and a complete (and correct) proof?
    \item[] Answer: \answerYes{}
    \item[] Justification: Main assumptions and proof sketches are in Sections 3--4. Theorem~1 (exact ODE solution) and Theorem~2 (degradation) are proved in Appendix~A. Propositions~1--3 (structural theorems) are proved in Appendices~B--D. Full broadband derivations and non-asymptotic bounds are in Appendices~E--F.

\item {\bf Experimental result reproducibility}
    \item[] Question: Does the paper fully disclose all the information needed to reproduce the main experimental results of the paper to the extent that it affects the main claims and/or conclusions of the paper (regardless of whether the code and data are provided or not)?
    \item[] Answer: \answerYes{}
    \item[] Justification: The protocol details, fairness controls, and reproducibility-critical settings are documented in Experiments and Appendix H.

\item {\bf Open access to data and code}
    \item[] Question: Does the paper provide open access to the data and code, with sufficient instructions to faithfully reproduce the main experimental results, as described in supplemental material?
    \item[] Answer: \answerNo{}
    \item[] Justification: At anonymous submission time we provide reproducibility instructions and provenance in Appendix H, but do not provide a fully open anonymized artifact bundle in this draft.

\item {\bf Experimental setting/details}
    \item[] Question: Does the paper specify all the training and test details (e.g., data splits, hyperparameters, how they were chosen, type of optimizer, etc.) necessary to understand the results?
    \item[] Answer: \answerYes{}
    \item[] Justification: Training budget, evaluation suite, fairness controls, and analysis protocol are described in Section 5 and Appendix H.

\item {\bf Experiment statistical significance}
    \item[] Question: Does the paper report error bars suitably and correctly defined or other appropriate information about the statistical significance of the experiments?
    \item[] Answer: \answerYes{}
    \item[] Justification: Multi-seed uncertainty is reported for 50M runs, and task-level paired analysis (bootstrap CI + exact sign-flip p-values) is reported for 8B LongBench.

\item {\bf Experiments compute resources}
    \item[] Question: For each experiment, does the paper provide sufficient information on the computer resources (type of compute workers, memory, time of execution) needed to reproduce the experiments?
    \item[] Answer: \answerYes{}
    \item[] Justification: Appendix H reports the fixed 8B budget and protocol settings, and the experimental setup sections provide the run-scale information used for reproducing reported comparisons.

\item {\bf Code of ethics}
    \item[] Question: Does the research conducted in the paper conform, in every respect, with the NeurIPS Code of Ethics \url{https://neurips.cc/public/EthicsGuidelines}?
    \item[] Answer: \answerYes{}
    \item[] Justification: This work is methodological analysis and controlled model evaluation; we are not aware of any deviations from the NeurIPS Code of Ethics.

\item {\bf Broader impacts}
    \item[] Question: Does the paper discuss both potential positive societal impacts and negative societal impacts of the work performed?
    \item[] Answer: \answerYes{}
    \item[] Justification: A dedicated Broader Impact section discusses potential benefits, misuse risks, and deployment safeguards for long-context systems.

\item {\bf Safeguards}
    \item[] Question: Does the paper describe safeguards that have been put in place for responsible release of data or models that have a high risk for misuse (e.g., pretrained language models, image generators, or scraped datasets)?
    \item[] Answer: \answerNA{}
    \item[] Justification: This submission does not release new high-risk models or scraped datasets as part of the paper artifact.

\item {\bf Licenses for existing assets}
    \item[] Question: Are the creators or original owners of assets (e.g., code, data, models), used in the paper, properly credited and are the license and terms of use explicitly mentioned and properly respected?
    \item[] Answer: \answerNo{}
    \item[] Justification: Existing assets are cited in text and references, but explicit per-asset license/terms tables are not fully enumerated in this draft.

\item {\bf New assets}
    \item[] Question: Are new assets introduced in the paper well documented and is the documentation provided alongside the assets?
    \item[] Answer: \answerNA{}
    \item[] Justification: The submission does not claim release of a new dataset/model benchmark asset.

\item {\bf Crowdsourcing and research with human subjects}
    \item[] Question: For crowdsourcing experiments and research with human subjects, does the paper include the full text of instructions given to participants and screenshots, if applicable, as well as details about compensation (if any)?
    \item[] Answer: \answerNA{}
    \item[] Justification: The experiments do not involve crowdsourcing or human-subject data collection.

\item {\bf Institutional review board (IRB) approvals or equivalent for research with human subjects}
    \item[] Question: Does the paper describe potential risks incurred by study participants, whether such risks were disclosed to the subjects, and whether Institutional Review Board (IRB) approvals (or an equivalent approval/review based on the requirements of your country or institution) were obtained?
    \item[] Answer: \answerNA{}
    \item[] Justification: No human-subject study is conducted in this work.

\item {\bf Declaration of LLM usage}
    \item[] Question: Does the paper describe the usage of LLMs if it is an important, original, or non-standard component of the core methods in this research? Note that if the LLM is used only for writing, editing, or formatting purposes and does not impact the core methodology, scientific rigorousness, or originality of the research, declaration is not required.
    \item[] Answer: \answerNA{}
    \item[] Justification: LLMs are not used as a non-standard component of the core method; the contribution is a mathematical framework plus controlled empirical evaluation.

\end{enumerate}

\appendix

\section{Proofs of Theorems 1 and 2}
\label{app:main_theorems}

\textbf{Proof of Theorem~1 (exact solution of the governing ODE).}
The joint functional after scale separation is
\begin{equation}
\mathcal{J}[\rho]=\frac{\alpha}{2}\int_0^1\rho(\phi)^2\,d\phi + \frac{\beta}{2}\int_0^1\!\!\int_0^1\rho(\phi_1)\rho(\phi_2)\min(\phi_1,\phi_2)\,d\phi_1\,d\phi_2 - \mu_F\int_0^1\rho(\phi)b^{-2\phi}\,d\phi + \lambda\!\left(\int_0^1\rho(\phi)\,d\phi-1\right).
\end{equation}
Setting the first variation to zero yields
\begin{equation}
\alpha\rho(\phi)+\beta\int_0^1\rho(\psi)\min(\phi,\psi)\,d\psi - \mu_F b^{-2\phi}+\lambda=0.
\end{equation}
Differentiating once in $\phi$,
$\alpha\rho'(\phi)+\beta\int_\phi^1\rho(\psi)\,d\psi - \mu_F(-2\ln b)b^{-2\phi}=0$,
which gives the boundary condition $\rho'(1)=-\frac{2\mu_F\ln b}{\alpha}b^{-2}$.
Differentiating again yields the non-homogeneous ODE
\begin{equation}
\alpha\rho''(\phi)-\beta\rho(\phi)=-\mu_F(2\ln b)^2 b^{-2\phi}.
\end{equation}
With $\tau=\sqrt{\beta/\alpha}$ and $\gamma_F=\mu_F(2\ln b)^2/\alpha$, this becomes $\rho''-\tau^2\rho=\gamma_F b^{-2\phi}$.
The homogeneous part has solutions $\cosh(\tau\phi)$, $\sinh(\tau\phi)$. A particular solution of the form $Pb^{-2\phi}$ gives $P=\gamma_F/(4(\ln b)^2-\tau^2)$, valid when $\tau\neq 2\ln b$.
Hence the general solution is $\rho^*(\phi)=C_1\cosh(\tau\phi)+C_2\sinh(\tau\phi)+\frac{\gamma_F}{4(\ln b)^2-\tau^2}b^{-2\phi}$,
with $C_1,C_2$ determined by $\rho'(1)$ and $\int_0^1\rho\,d\phi=1$. \hfill $\square$

\medskip
\textbf{Proof of Theorem~2 (asymptotic degradation of EVQ).}
The EVQ warp is $\phi_k(\tau)=1-\frac{1}{\tau}\operatorname{arcsinh}\!\bigl((1-u_k)\sinh\tau\bigr)$ with $u_k=(k+0.5)/N$.
As $\tau\to 0$, $\sinh\tau\to\tau+\mathcal{O}(\tau^3)$ and $\operatorname{arcsinh}(x)\to x+\mathcal{O}(x^3)$ for small $x$. Hence
\begin{equation}
\phi_k(\tau)=1-\frac{1}{\tau}\operatorname{arcsinh}\!\bigl((1-u_k)\tau+\mathcal{O}(\tau^3)\bigr)
=1-\frac{(1-u_k)\tau+\mathcal{O}(\tau^3)}{\tau}=u_k+\mathcal{O}(\tau^2).
\end{equation}
But $u_k=(k+0.5)/N$ is precisely the midpoint quantisation of the uniform CDF on $[0,1]$, which corresponds to geometric (log-linear) RoPE spacing. \hfill $\square$

\section{Proof of Proposition 1 (uniform prior)}
\label{app:prop1}
We prove that a uniform distance prior yields asymptotically constant diagonal energy.

\textbf{Setup.} Let $D(\Delta)=\frac{1}{L}\mathbb{1}_{\Delta\in[0,L]}$. Then
\begin{equation}
\Db(\xi)=\int_0^L \frac{1}{L}\cos(\xi\Delta)\,d\Delta=\frac{\sin(\xi L)}{\xi L}=\sinc(\xi L).
\label{eq:Db_sinc}
\end{equation}
Hence $\Ediag(\phi)=\tfrac{1}{2}\bigl(1+\sinc(2b^{-\phi}L)\bigr)$.

\textbf{Uniform decay as $L\to\infty$.} Fix $b$ and $\phi\in[0,1]$. Then $\xi(\phi)=2b^{-\phi}\ge 2/b>0$. As $L\to\infty$, $\sinc(\xi(\phi)L)\to 0$ and $|\sinc(\xi(\phi)L)|\le 1/(\xi(\phi)L)$. In particular, for any $\phi\in[0,1]$,
\begin{equation}
\left|\Ediag(\phi)-\frac{1}{2}\right|=\frac{1}{2}|\sinc(\xi(\phi)L)|\le \frac{1}{2\xi(\phi)L}\le \frac{b}{4L}\to 0,
\label{eq:Ediag_uniform}
\end{equation}
which implies uniform convergence $\Ediag(\phi)\to 1/2$.

\textbf{Conclusion.} By Lemma~1, $\rho^{\star}(\phi)\propto 1/\Ediag(\phi)$ converges to a constant density, corresponding to geometric RoPE allocation. \hfill $\square$

\section{Proof of Proposition 2 (power-law prior, corrected Ci asymptotics)}
\label{app:prop2}
This appendix corrects a common algebraic pitfall: using $\Ci(z)\approx \gamma+\ln z$ outside its validity regime. The expansion $\Ci(z)=\gamma+\ln z+\mathcal{O}(z^2)$ holds for $z\ll 1$; for $z\gg 1$, $\Ci(z)$ decays oscillatory as $\Ci(z)=\frac{\sin z}{z}+\mathcal{O}(1/z^2)$.

\textbf{Power-law prior and diagonal energy.} Let $D(\Delta)=\frac{1}{\Delta\ln L}\mathbb{1}_{\Delta\in[1,L]}$. Then
\begin{equation}
\Db(\xi)=\frac{1}{\ln L}\int_1^L \frac{\cos(\xi\Delta)}{\Delta}\,d\Delta
=\frac{1}{\ln L}\int_{\xi}^{\xi L} \frac{\cos u}{u}\,du
=\frac{\Ci(\xi L)-\Ci(\xi)}{\ln L}.
\label{eq:Db_Ci}
\end{equation}
Hence
\begin{equation}
\Ediag(\phi)=\frac{1}{2}\left(1+\frac{\Ci(2b^{-\phi}L)-\Ci(2b^{-\phi})}{\ln L}\right).
\label{eq:Ediag_Ci_exact}
\end{equation}

\textbf{Bulk asymptotics and sign (fixing the algebraic error).} Let $\xi=2b^{-\phi}$. For $\xi$ (bulk region, $\phi$ not too close to 0 so that $\xi\ll 1$),
\begin{equation}
\Ci(\xi)\approx \gamma+\ln\xi=\gamma+\ln2-\phi\ln b.
\label{eq:Ci_small}
\end{equation}
For $\xi L=2b^{-\phi}L$ (when $L\gg b$, $\xi L\gg 1$),
\begin{equation}
\Ci(\xi L)\approx \frac{\sin(\xi L)}{\xi L}\approx 0\qquad\text{(oscillatory decay)}.
\label{eq:Ci_large}
\end{equation}
Therefore,
\begin{equation}
\Db(2b^{-\phi})\approx \frac{0-(\gamma+\ln2-\phi\ln b)}{\ln L}=\frac{\phi\ln b-\gamma-\ln2}{\ln L},
\label{eq:Db_affine}
\end{equation}
and
\begin{equation}
\Ediag(\phi)\approx \frac{1}{2}\left(1+\frac{\phi\ln b-\gamma-\ln2}{\ln L}\right),
\label{eq:Ediag_affine_app}
\end{equation}
which is strictly increasing in $\phi$. Therefore $\rho^{\star}(\phi)\propto 1/\Ediag(\phi)$ is strictly decreasing: a high-frequency bias.

\textbf{Amplitude remains bounded.} Using the affine bulk form, one obtains an end-to-end ratio
\begin{equation}
\frac{\rho^{\star}(0)}{\rho^{\star}(1)}=\frac{\Ediag(1)}{\Ediag(0)}
\approx
\frac{\ln L+\ln b-\gamma-\ln2}{\ln L-\gamma-\ln2}
=1+\frac{\ln b}{\ln L-\gamma-\ln2},
\label{eq:amplitude_ratio}
\end{equation}
which is $1+\mathcal{O}(\ln b/\ln L)$ when $\ln L\gg \ln b$.
This expression is a bulk proxy: at $\phi=0$ (so $\xi=2$), one should use the exact constant $\Ci(2)$ instead of $\gamma+\ln2$. The boundary correction remains finite and does not alter the bounded-amplitude conclusion.

\textbf{Takeaway.} The incorrect step to avoid is replacing both $\Ci(\xi)$ and $\Ci(\xi L)$ by $\gamma+\ln(\cdot)$, which would cancel to $\ln L$ and erase $\phi$-dependence. The correct regime split is: $\Ci(\xi)\approx \gamma+\ln\xi$ for $\xi\ll 1$ but $\Ci(\xi L)\approx 0$ for $\xi L\gg 1$. \hfill $\square$

\section{Proof of Proposition 3 (resonance condition and deep local minima)}
\label{app:prop3}
Let $D(\Delta)=\lambda\delta(\Delta-s)+(1-\lambda)\delta(\Delta-\ell)$. Then
\begin{equation}
\Ediag(\phi)=\frac{1}{2}\Bigl(1+\lambda\cos(2b^{-\phi}s)+(1-\lambda)\cos(2b^{-\phi}\ell)\Bigr).
\label{eq:Ediag_bimodal_app}
\end{equation}

\textbf{Resonance bound.} Assume there exist $\phi_\star\in[0,1]$, integers $m,n$, and nonnegative errors $\varepsilon_s,\varepsilon_\ell$ such that
\begin{equation}
\left|2b^{-\phi_\star}s-(2m+1)\pi\right|\le \varepsilon_s,\qquad
\left|2b^{-\phi_\star}\ell-(2n+1)\pi\right|\le \varepsilon_\ell.
\label{eq:resonance_cond_app}
\end{equation}
Write
\begin{equation}
2b^{-\phi_\star}s=(2m+1)\pi+\delta_s,\quad
2b^{-\phi_\star}\ell=(2n+1)\pi+\delta_\ell,\quad
|\delta_s|\le\varepsilon_s,\;|\delta_\ell|\le\varepsilon_\ell.
\end{equation}
Then
\begin{equation}
\cos(2b^{-\phi_\star}s)=-\cos\delta_s,\qquad
\cos(2b^{-\phi_\star}\ell)=-\cos\delta_\ell,
\end{equation}
and therefore
\begin{equation}
\Ediag(\phi_\star)=\frac{1}{2}\Bigl(1-\lambda\cos\delta_s-(1-\lambda)\cos\delta_\ell\Bigr)
\le
\frac{1}{2}\Bigl(1-\lambda\cos\varepsilon_s-(1-\lambda)\cos\varepsilon_\ell\Bigr).
\label{eq:resonance_upper_app}
\end{equation}
Using $\cos x\ge 1-x^2/2$ for small $x$ gives
\begin{equation}
\Ediag(\phi_\star)\le \frac{\lambda\varepsilon_s^2+(1-\lambda)\varepsilon_\ell^2}{4}
=\mathcal{O}(\varepsilon_s^2+\varepsilon_\ell^2).
\label{eq:resonance_small_app}
\end{equation}
Hence near-resonant alignments create arbitrarily deep diagonal minima.

\textbf{Consequence for the inverse law.} By Lemma~1, $\rho^{\star}(\phi)\propto 1/\Ediag(\phi)$ spikes near such minima, which is brittle and can diverge from the optimizer of the full functional where off-diagonal coupling smooths the landscape. \hfill $\square$
For fixed $b$, larger long-range distances $(s,\ell)$ increase phase oscillation counts over $\phi\in[0,1]$, heuristically increasing the chance of near-resonant points.

\section{Diagonal approximation and residual analysis (summary)}
Expanding the kernel,
\begin{equation}
K(\phi_1,\phi_2)=\int D(\Delta)\cos(\omega(\phi_1)\Delta)\cos(\omega(\phi_2)\Delta)\,d\Delta.
\label{eq:kernel_expand}
\end{equation}
Using $\cos a\cos b=\tfrac{1}{2}(\cos(a-b)+\cos(a+b))$ shows that $K(\phi,\phi)$ contains a constant $1/2$ plus an oscillatory term involving $\Db(\omega(\phi))$, yielding the diagonal surrogate. Off-diagonal terms involve $\Db(\omega(\phi_1)\pm\omega(\phi_2))$ and are small when $D$ is sufficiently smooth and broadband. A representative finite-base audit at $b=10^4$ reports an $\mathcal{O}(1/\ln b)$ residual of about $11\%$.

\section{Structural stability beyond the diagonal}
\subsection{Brownian covariance}
In the broadband limit $b\to\infty$ under a power-law prior, the off-diagonal kernel converges to
\begin{equation}
K_{\mathrm{off}}(\phi_1,\phi_2)\to \frac{1}{2}\min(\phi_1,\phi_2),
\label{eq:Koff_brownian}
\end{equation}
the covariance of Brownian motion. We use the regularized functional
\begin{equation}
\mathcal{J}[\rho]
=\frac{1}{2}\int_0^1 \rho(\phi)^2\,d\phi
+\frac{1}{2}\int_0^1\!\!\int_0^1 \rho(\phi_1)\rho(\phi_2)\min(\phi_1,\phi_2)\,d\phi_1d\phi_2
+\mu\!\left(\int_0^1 \rho(\phi)\,d\phi-1\right).
\label{eq:J_brownian_app}
\end{equation}
The first variation yields
\begin{equation}
\rho(\phi)+\int_0^1 \rho(\psi)\min(\phi,\psi)\,d\psi+\mu=0.
\label{eq:first_var_app}
\end{equation}
Differentiating once gives
\begin{equation}
\rho'(\phi)+\int_{\phi}^{1}\rho(\psi)\,d\psi=0,
\end{equation}
and differentiating again gives
\begin{equation}
\rho''(\phi)-\rho(\phi)=0,\qquad \rho'(1)=0.
\label{eq:cosh_ode_app}
\end{equation}
Hence $\rho\propto\cosh(1-\phi)$. If one keeps only the pure off-diagonal term in \eqref{eq:J_brownian_app}, the resulting equation is different ($\rho''=0$), so the diagonal regularizer is essential to recover the observed bounded profile.

\subsubsection{Uniform validity}
At $\phi=0$ directly evaluate the exact off-diagonal kernel:
\begin{equation}
K_{\mathrm{off}}(0,\psi)=-\frac{1}{4\ln b}\Bigl[\Ci(1-\omega_2)+\Ci(1+\omega_2)\Bigr],
\label{eq:Koff_phi0}
\end{equation}
where $\omega_2=\omega(\psi)=b^{-\psi}\in(0,1)$. Using a symmetric Taylor expansion around argument $1$, $\Ci(1\pm\omega_2)=\Ci(1)\pm \Ci'(1)\omega_2+\mathcal{O}(\omega_2^2)$, the first-order terms cancel, yielding the pointwise boundary correction
\begin{equation}
\rho_{\mathrm{exact}}(0)-\rho_{\mathrm{out}}(0)=\frac{\Ci(1)}{2\ln b}\approx 0.018\qquad (b=10{,}000).
\label{eq:boundary_correction}
\end{equation}
Since $\cosh(1)\approx 1.54$, the discrepancy is only $\sim 1.2\%$. There is no boundary singularity.

\subsection{Non-commutativity}
Under a uniform prior, the physically correct limit is $L\to\infty$ with finite $d$; reversing the order ($d\to\infty$ first) introduces a Jacobian distortion that changes the inferred density. This explains why some continuous approximations can contradict discrete behavior; see the discussion in Remark~1.

\subsection{Fredholm correction}
At finite base, the full optimization yields a Fredholm integral equation. We expand
\begin{equation}
\rho=\rho_0+\varepsilon\rho_1,\qquad \varepsilon=1/\ln b,\qquad \rho_0=\cosh(1-\phi),
\end{equation}
and obtain the first-order correction
\begin{equation}
\rho_1(\phi)=-\int_0^1 \min(\phi,\psi)\cosh(1-\psi)\,d\psi.
\label{eq:rho1_int}
\end{equation}
This admits the closed form
\begin{equation}
\rho_1(\phi)=\cosh(1-\phi)-\cosh(1),
\label{eq:rho1_closed}
\end{equation}
which is strictly negative for all $\phi>0$. Intuitively, it acts as a uniform flattening penalty away from the high-frequency endpoint. The corrected solution is
\begin{equation}
\rho(\phi)\approx (1+\varepsilon)\cosh(1-\phi)-\varepsilon\cosh(1).
\label{eq:rho_corrected}
\end{equation}

\textbf{Reconciling remark.}
The $b$-independence of the leading $\cosh(1-\phi)$ profile reflects the $b\to\infty$ broadband limit, where diagonal variation is suppressed. The Fredholm correction reintroduces finite-base dependence through $\varepsilon=1/\ln b$. Thus Theorem~2 (diagonal, finite $b$) and the cosh solution (full, broadband limit) describe complementary regimes rather than conflicting predictions.

\subsection{Waterbed inequality}
Define $\mathcal{I}(\phi)=c\rho(\phi)b^{-2\phi}$ and note the proxy relation $E(\phi)\ge 1/\mathcal{I}(\phi)$. By Jensen's inequality,
\begin{equation}
\int_0^1 \ln E(\phi)\,d\phi \ge \ln b-\ln c.
\label{eq:waterbed_app}
\end{equation}
This is an information-theoretic scaling bound: it formalizes an integrated trade-off between phase-collision suppression and local resolution, without claiming a pointwise tight operational bound.

\section{Additional non-asymptotic and discretization guarantees}

\textbf{Proposition G.1 (Broadband decomposition and structural invariance).}
\textit{Let $D(\Delta)=\frac{1}{\Delta\ln L}\mathbb{1}_{\Delta\in[1,L]}$. In the broadband limit, the full kernel decomposes into a diagonal singular ridge plus a smooth off-diagonal Brownian covariance. Modeling the diagonal ridge with a Tikhonov weight $\alpha>0$, consider}
\begin{equation}
\mathcal{J}_{\alpha}[\rho]
=\frac{\alpha}{2}\!\int_0^1\!\rho(\phi)^2\,d\phi
+\frac{1}{2}\!\int_0^1\!\!\int_0^1\!\rho(\phi_1)\rho(\phi_2)\min(\phi_1,\phi_2)\,d\phi_1d\phi_2
+\mu\!\left(\int_0^1\!\rho(\phi)\,d\phi-1\right).
\label{eq:J_alpha}
\end{equation}
\textit{Its optimizer is unique (up to normalization) and satisfies}
\begin{equation}
\rho_{\alpha}^{\star}(\phi)\propto \cosh\!\left(\frac{1-\phi}{\sqrt{\alpha}}\right),
\label{eq:rho_alpha_cosh}
\end{equation}
\textit{hence is strictly decreasing and strictly convex on $[0,1)$.}

\textbf{Proof.}
For $\phi_1\neq\phi_2$ and (without loss of generality) $\phi_1>\phi_2$, we have $\omega_1\ll\omega_2$ and therefore $\omega_1\pm\omega_2=\pm\omega_2+o(1)$. Using $\Db(\xi)=\frac{\Ci(\xi L)-\Ci(\xi)}{\ln L}$ with $\Ci(\xi L)\approx 0$ and $\Ci(\xi)\approx \gamma+\ln\xi$ in the broadband asymptotic regime yields an off-diagonal term affine in the smaller log-frequency coordinate, i.e., proportional to $\min(\phi_1,\phi_2)$. At $\phi_1=\phi_2$, $\Db(0)=1$ creates a diagonal ridge contribution; in the continuous model we keep that contribution via the local $L_2$ term with weight $\alpha$.

The first variation of \eqref{eq:J_alpha} gives
\begin{equation}
\alpha\rho(\phi)+\int_0^1 \rho(\psi)\min(\phi,\psi)\,d\psi+\mu=0.
\label{eq:Jalpha_stationary}
\end{equation}
Differentiating once,
\begin{equation}
\alpha\rho'(\phi)+\int_{\phi}^{1}\rho(\psi)\,d\psi=0,
\label{eq:Jalpha_first_deriv}
\end{equation}
which implies $\rho'(1)=0$. Differentiating again yields
\begin{equation}
\alpha\rho''(\phi)-\rho(\phi)=0,\qquad \rho'(1)=0.
\label{eq:Jalpha_ode}
\end{equation}
Thus $\rho(\phi)=A\cosh((1-\phi)/\sqrt{\alpha})$, and normalization fixes $A$. Since $\sinh(x)>0$ and $\cosh(x)>0$ for $x>0$, we have $\rho'_{\alpha}(\phi)<0$ and $\rho''_{\alpha}(\phi)>0$ on $[0,1)$. Strict convexity of $\mathcal{J}_{\alpha}$ for $\alpha>0$ gives uniqueness. \hfill $\square$
Physically, $\alpha$ controls the relative strength of the diagonal ridge; for finite discrete heads, this term reflects the non-vanishing diagonal mass contribution (order $1/N$). In the formal limit $\alpha\to0^+$, the model degenerates to the pure Brownian regime with linear stationarity, while any $\alpha>0$ preserves the observed monotone-convex high-frequency bias.

\textbf{Proposition G.2 (Finite-$(b,L)$ residual bound for the affine surrogate).}
\textit{Let}
\begin{equation}
\widetilde{E}_{\mathrm{diag}}(\phi)
=\frac{1}{2}\left(1+\frac{\phi\ln b-\gamma-\ln2}{\ln L}\right),
\label{eq:Ediag_tilde}
\end{equation}
\textit{and define $\Delta E(\phi)=|E_{\mathrm{diag}}(\phi)-\widetilde{E}_{\mathrm{diag}}(\phi)|$. For finite $b,L$:}
\begin{equation}
\Delta E(\phi)\le
\frac{1}{2\ln L}\left(b^{-2\phi}+\frac{b^{\phi}}{L}\right).
\label{eq:finite_residual_bound}
\end{equation}

\textbf{Proof.}
From \eqref{eq:Ediag_Ci_exact},
\begin{equation}
E_{\mathrm{diag}}(\phi)=\frac{1}{2}\left(1+\frac{\Ci(2b^{-\phi}L)-\Ci(2b^{-\phi})}{\ln L}\right).
\end{equation}
For $x>0$, using $\Ci(x)=\gamma+\ln x-\int_0^x\frac{1-\cos t}{t}\,dt$ and $1-\cos t\le t^2/2$,
\begin{equation}
\left|\Ci(x)-(\gamma+\ln x)\right|
\le \int_0^x \frac{t}{2}\,dt=\frac{x^2}{4}.
\end{equation}
At $x=2b^{-\phi}$, this gives a near-field error $\le b^{-2\phi}$. For $X>0$, integration by parts implies
\begin{equation}
|\Ci(X)|=\left|\frac{\sin X}{X}-\int_X^\infty\frac{\sin t}{t^2}\,dt\right|\le \frac{2}{X},
\end{equation}
so for $X=2b^{-\phi}L$ the far-field error is $\le b^{\phi}/L$. Combining by triangle inequality and dividing by $2\ln L$ yields \eqref{eq:finite_residual_bound}. \hfill $\square$
For the two finite-base settings used in Section~4.3 ($L=16384$ with $b=10^4$ and $b=5\times10^5$), this bound is consistent with the observed high-$R^2_{\mathrm{mid}}$ calibration and small mid-band mean absolute residuals.

\textbf{Proposition G.3 (Regularized inverse law and concentration collapse).}
\textit{Assume a resonant trap at $\phi_{\star}$ with local behavior $E_{\mathrm{diag}}(\phi)\approx c(\phi-\phi_{\star})^2$, $c>0$. Define}
\begin{equation}
\rho_{\epsilon}^{\star}(\phi)=\frac{1}{Z_{\epsilon}}\frac{1}{E_{\mathrm{diag}}(\phi)+\epsilon},\qquad \epsilon>0.
\label{eq:rho_eps_def}
\end{equation}
\textit{Then $\|\rho_{\epsilon}^{\star}\|_{\infty}=\mathcal{O}(1/\sqrt{\epsilon})$. As $\epsilon\to0^+$, $\rho_{\epsilon}^{\star}$ converges weakly to $\delta(\phi-\phi_{\star})$.}

\textbf{Proof.}
Near $\phi_{\star}$ we approximate
\begin{equation}
Z_{\epsilon}\approx \int_{-\infty}^{\infty}\frac{d\phi}{c(\phi-\phi_{\star})^2+\epsilon}
=\frac{\pi}{\sqrt{c\epsilon}}.
\label{eq:Zeps_eval}
\end{equation}
Hence
\begin{equation}
\rho_{\epsilon}^{\star}(\phi_{\star})
=\frac{1}{\epsilon Z_{\epsilon}}
=\frac{\sqrt{c}}{\pi\sqrt{\epsilon}}
=\mathcal{O}\!\left(\frac{1}{\sqrt{\epsilon}}\right).
\end{equation}
For any continuous test function $f$, the normalized Cauchy kernel in \eqref{eq:rho_eps_def} satisfies
\begin{equation}
\lim_{\epsilon\to0^+}\int_0^1 f(\phi)\rho_{\epsilon}^{\star}(\phi)\,d\phi=f(\phi_{\star}),
\end{equation}
which is weak convergence to $\delta(\phi-\phi_{\star})$. \hfill $\square$
In sparse discrete settings (e.g., $d=128$ and $N=64$ frequencies), this collapse appears as density concentrating on very few bins, effectively starving the remaining frequency coordinates.

\textbf{Proposition G.4 (Pointwise off-diagonal decay for distinct frequencies).}
\textit{Under $D(\Delta)=\frac{1}{L}\mathbb{1}_{\Delta\in[0,L]}$, let}
\begin{equation}
K_{\mathrm{off}}(\phi_1,\phi_2)
=\frac{1}{2}\left[\sinc((\omega_1-\omega_2)L)+\sinc((\omega_1+\omega_2)L)\right],\quad \omega_i=b^{-\phi_i}.
\end{equation}
\textit{For any pair with $\phi_1\neq\phi_2$, $K_{\mathrm{off}}(\phi_1,\phi_2)=\mathcal{O}(1/L)$ as $L\to\infty$. Moreover, on any separated set $\{|\omega_1-\omega_2|\ge \delta>0\}$, the bound is uniform.}

\textbf{Proof.}
Using $|\sinc(x)|\le 1/|x|$,
\begin{equation}
|K_{\mathrm{off}}(\phi_1,\phi_2)|
\le \frac{1}{2}\left(\frac{1}{|\omega_1-\omega_2|L}+\frac{1}{(\omega_1+\omega_2)L}\right)
=\mathcal{O}\!\left(\frac{1}{L}\right),
\end{equation}
for fixed $(\phi_1,\phi_2)$ with $\omega_1\neq\omega_2$. If $|\omega_1-\omega_2|\ge\delta$, then
\begin{equation}
|K_{\mathrm{off}}(\phi_1,\phi_2)|\le \frac{1}{2L}\left(\frac{1}{\delta}+\frac{1}{\omega_1+\omega_2}\right)\le \frac{C_{\delta}}{L},
\end{equation}
which is uniform on the separated set. \hfill $\square$
Under the power-law prior, an analogous asymptotic argument gives logarithmic attenuation with $L$ for corresponding off-diagonal terms.

\textbf{Proposition G.5 (Quantization error and anchoring necessity).}
\textit{Let $\phi_k=F^{-1}\!\left(\frac{k+0.5}{N}\right)$ with $F(\phi)=\int_0^\phi \rho(t)\,dt$, and define $g(u)=E_{\mathrm{diag}}(F^{-1}(u))$ on $u\in[0,1]$. Then midpoint quantization gives}
\begin{equation}
\left|\frac{1}{N}\sum_{k=0}^{N-1} g\!\left(\frac{k+0.5}{N}\right)-\int_0^1 g(u)\,du\right|
\le \frac{1}{24N^2}\sup_{u\in[0,1]}|g''(u)|.
\label{eq:midpoint_error}
\end{equation}
\textit{Moreover,}
\begin{equation}
g''(u)
=\frac{E_{\mathrm{diag}}''(\phi)}{\rho(\phi)^2}
-\frac{E_{\mathrm{diag}}'(\phi)\rho'(\phi)}{\rho(\phi)^3},
\label{eq:gpp_expression}
\end{equation}
\textit{hence}
\begin{equation}
\mathcal{E}_{\mathrm{quant}}
\le
\frac{1}{24N^2}
\sup_{\phi\in[0,1]}
\left|
\frac{E_{\mathrm{diag}}''(\phi)}{\rho(\phi)^2}
-\frac{E_{\mathrm{diag}}'(\phi)\rho'(\phi)}{\rho(\phi)^3}
\right|.
\label{eq:quant_bound_final}
\end{equation}
\textit{If $\inf_{\phi}\rho(\phi)\to0$, the bound can blow up; enforcing $\inf_{\phi}\rho(\phi)\ge c_0>0$ through anchoring prevents this singular behavior.}

\textbf{Proof.}
With $u=F(\phi)$ and $d\phi/du=1/\rho(\phi)$, midpoint quadrature on $g$ yields \eqref{eq:midpoint_error}. By chain rule,
\begin{equation}
g'(u)=E_{\mathrm{diag}}'(\phi)\frac{d\phi}{du}
=\frac{E_{\mathrm{diag}}'(\phi)}{\rho(\phi)},
\end{equation}
and differentiating again gives \eqref{eq:gpp_expression}. Substituting into \eqref{eq:midpoint_error} yields \eqref{eq:quant_bound_final}. The $1/\rho^3$ term shows why near-zero density creates discretization instability, while a positive density floor prevents blow-up. \hfill $\square$

\section{Experimental Hyperparameter Configurations}
This appendix summarizes the reproducibility-critical settings used for all reported results in a submission-ready format.

\paragraph{H.1 From-scratch TinyStories runs.}
\begin{table}[H]
\centering
\begin{tabular}{p{0.28\linewidth}p{0.66\linewidth}}
\toprule
Item & Configuration \\
\midrule
Models & 50M (3 seeds), 100M (single seed), 350M (single seed) \\
Dataset & TinyStories~\cite{eldan2023tinystories} \\
Objective & Autoregressive next-token prediction (PPL reported) \\
Train/eval lengths & Train at 2K; evaluate at 2K and 16K (plus intermediates where available) \\
Compared schedules & Geometric ($\tau=0$) vs.\ EVQ ($\tau>0$) (same optimizer/data path per scale) \\
Reported metric & Perplexity at fixed context lengths \\
\bottomrule
\end{tabular}
\caption{Core settings for Table~\ref{tab:tinystories_scaling}.}
\label{tab:app_h_tinystories}
\end{table}

\paragraph{H.2 50M YaRN comparison under matched budget.}
\begin{table}[H]
\centering
\begin{tabular}{p{0.28\linewidth}p{0.66\linewidth}}
\toprule
Item & Configuration \\
\midrule
Model scale & 50M \\
Methods & Geometric ($\tau=0$), YaRN (progressive), EVQ ($\tau>0$) \\
Budget control & Matched training budget and data pipeline across methods \\
Evaluation lengths & 2K, 4K, 8K, 12K, 16K \\
Reported metric & Perplexity by length \\
\bottomrule
\end{tabular}
\caption{Core settings for the controlled YaRN comparison.}
\label{tab:app_h_yarn}
\end{table}

\paragraph{H.3 8B controlled protocol.}
\begin{table}[H]
\centering
\begin{tabular}{p{0.28\linewidth}p{0.66\linewidth}}
\toprule
Item & Configuration \\
\midrule
Base model & Llama-3-8B~\cite{grattafiori2024llama3} \\
Adaptation regime & LoRA, 600 steps, context $L=16\mathrm{K}$ \\
Alignment setting & Minimal adaptation (no instruction-format alignment; no large QA fine-tuning) \\
Implementation control & Shared \texttt{inv\_freq.copy()} injection path; no built-in \texttt{rope\_scaling} API differences \\
Compared schedules & Geometric ($\tau=0$), PI, YaRN, Diagonal-only, EVQ ($\tau>0$) \\
Evaluation suite & NIAH@16K~\cite{kamradt2023niah}, LongBench (6 tasks)~\cite{bai2023longbench}, Passkey-TF margin@16K \\
Statistics policy & Task-level paired analysis ($n=6$); full per-sample traces not exported in this archived package \\
\bottomrule
\end{tabular}
\caption{Core settings for Table~\ref{tab:8b_fair}.}
\label{tab:app_h_8b}
\end{table}


\paragraph{H.4 Qwen-2.5-7B dual-seed controlled protocol.}
\begin{table}[H]
\centering
\begin{tabular}{p{0.28\linewidth}p{0.66\linewidth}}
\toprule
Item & Configuration \\
\midrule
Base model & Qwen-2.5-7B-Instruct \\
Adaptation regime & LoRA, 400 steps, batch=2, seq\_len=16K, lr=2e-5 \\
Training data & WikiText-style corpora (same pipeline as 8B suite) \\
Seeds & 42, 1337 \\
Implementation control & Shared \texttt{inv\_freq.copy()} injection path; EVQ with $\tau$ calibrated via anchor\_factor=4.0, slope=20.0, center\_ratio=0.70 \\
Compared schedules & Geometric baseline ($\tau=0$) vs.\ EVQ ($\tau>0$) \\
Evaluation suite & LongBench-21 (21 tasks, full coverage including single-doc QA, multi-doc QA, summarization, few-shot, synthetic, code) \\
Statistics policy & Per-sample bootstrap with FDR (Benjamini-Hochberg) correction; $n{=}4750$ samples per task pair; dual-seed replication \\
Quality checks & num\_selected = num\_scored = 4750; generation\_error = 0; parse\_fail = 0; template leakage $<1\%$ \\
\bottomrule
\end{tabular}
\caption{Core settings for Tables~\ref{tab:qwen_dual_seed} and~\ref{tab:qwen_task_family}.}
\label{tab:app_h_qwen}
\end{table}


\end{document}
